\documentclass[11pt]{article}
\usepackage{amsmath,amsthm,amssymb}
\usepackage[margin=1in]{geometry}
\usepackage{enumitem}
\usepackage{booktabs}
\usepackage{hyperref}

\newtheorem{theorem}{Theorem}
\newtheorem{lemma}[theorem]{Lemma}
\newtheorem{proposition}[theorem]{Proposition}
\newtheorem{corollary}[theorem]{Corollary}
\newtheorem{conjecture}[theorem]{Conjecture}
\theoremstyle{definition}
\newtheorem{definition}[theorem]{Definition}
\newtheorem{remark}[theorem]{Remark}
\newtheorem{example}[theorem]{Example}

\title{Lower bound program for $\lambda = (4, 3, 1^q)$:\\
  five-vector construction, A--B separation, and the\\
  position-of-letter contamination obstruction}
\author{Clio}
\date{13 May 2026 (afternoon prove session)}

\begin{document}
\maketitle

\begin{abstract}
The May-13 4-3-1n upper bound paper
(\texttt{2026-05-13-4-3-1n-upper-bound.tex}) proved
$\dim B^{(\lambda)}_{j_0} V_\lambda \le 5$ structurally at $q \ge 5$
for $\lambda = (4, 3, 1^q)$, with computational verification of dim
$= 5$ at $q \in \{3, 4\}$. The matching \emph{structural} lower bound
$\dim \ge 5$ was deferred. This paper records partial progress toward
that lower bound and identifies the precise technical obstruction.

\textbf{What we establish.} (i) An explicit five-vector construction:
two vectors from $\Phi_A^\lambda$ lifting a basis of $B^{(\mu_A)}$
($\mu_A = (3, 3, 1^{q-1})$, dim 2), three vectors from $\Phi_B^\lambda$
lifting a basis of $B^{(\mu_B)}$ ($\mu_B = (4, 2, 1^{q-1})$, dim 3).
(ii) Structural \emph{A vs B} separation via position-of-$n$: the
exclusive cells $c_1 = (1, 4)$ for $A$ and $c_2 = (2, 3)$ for $B$
distinguish the two corner-pair contributions at the $\lambda$ level,
and the position-of-$n$ projection commutes with the outer chain
$\mathcal{O}_\lambda$. (iii) An analysis showing that the natural
\emph{within-$X$} position-of-letter separators (analogous to May-12
evening-3's position-of-$(n-1)$ for $(4, 2, 1^*)$) do \emph{not}
propagate cleanly through the outer chain for either $X \in \{A, B\}$
of the $(4, 3, 1^q)$ shape.

\textbf{The obstruction.} For $(4, 2, 1^q)$ the outer chain has three
factors $(T_{n-4}+1)(T_{n-3}+1)(T_{n-2}+1)$, and the separator
position-of-$(n-1)$ is preserved by the first two factors, with only
the last factor producing controlled motion via a single
$T_{n-2}$-$2$-block. The cells reachable by this single move form a
\emph{$2$-element set} that has clean exclusive elements. For
$(4, 3, 1^q)$ the outer chain has \emph{four} factors
$(T_{n-5}+1) \cdots (T_{n-2}+1)$, the natural separator
position-of-$(n-2)$ is moved by both $(T_{n-3}+1)$ and $(T_{n-2}+1)$,
and the two-step motion enlarges the reachable cell set to include
cells originally exclusive to the other branch. The contamination is
not an artefact of the proof method but a structural fact: pos-of-
$(n-2)$ at the $\lambda$ level reaches the same set of cells starting
from either inner basis vector.

\textbf{Path forward.} We identify the \emph{tracer-pair} technique
(this session's memory note `tracer-pair-technique`) as the
appropriate next tool: rather than separating $F_A^1, F_A^2$ via a
position-of-letter projection, identify two specific $\lambda$-SYTs
$T^*_1, T^*_2$ and compute the seminormal coefficients
$\langle v_{T^*_i}, F_A^j\rangle$ as products of Hoefsmit factors;
exhibit a $2 \times 2$ tracer matrix with $\mathbb{Q}(q)^\times$
determinant. The chain of operations is finite and the computation
is explicit; we sketch its structure but do not complete it in this
paper.

\textbf{Status.} Partial. The A--B separation and the obstruction
analysis are fully structural. The within-$X$ linear independence
reduces to a tracer-pair Hoefsmit computation that we leave for a
follow-up paper.
\end{abstract}

\section{Setup}\label{sec:setup}

Conventions follow the May-13 4-3-1n upper bound paper
(\texttt{2026-05-13-4-3-1n-upper-bound.tex}, henceforth
\textbf{May-13-UB}). $H_q(S_n)$ is the Iwahori--Hecke algebra over
$\mathbb{Q}(q)$ with generators $T_1, \ldots, T_{n-1}$ obeying
$(T_i - q)(T_i + 1) = 0$ and the standard braid relations.
$V_\lambda$ is the seminormal cell module for $\lambda \vdash n$ in
asymmetric Murphy normalisation ($b' = 1$). $E_j^\pm$ denote the
$q$- and $(-1)$-eigenspaces of $T_j$. For $1 \le k \le n$,
\[
   R'_k \;:=\; (T_{k-1}{+}1)(T_{k-2}{+}1)\cdots(T_1{+}1),
   \qquad
   B^{(\lambda)}_j V_\lambda \;:=\; R'_j R'_{j+1}\cdots R'_n V_\lambda.
\]
Sharp index $j_0(\lambda) := \ell(\lambda) + 1$.

\paragraph{Throughout this paper.}
$\lambda := (4, 3, 1^q)$ with $q \ge 3$, $n := q + 7$, $\ell := q + 2$,
$j_0 := q + 3$. Removable corners
\[
   c_1 = (1, 4), \quad c_2 = (2, 3), \quad c_3 = c_* = (q+2, 1).
\]
$c_1, c_2$ are length-preserving; $c_3$ is the column-$1$ bottom.
Same-row pair at row $2$: cells $(2, 2), (2, 3)$.

\paragraph{Inner partitions.}
\begin{align*}
   \mu_A &:= \lambda \setminus \{c_1, c_3\} = (3, 3, 1^{q-1}),
      &|\mu_A| = q+5 \\
   \mu_B &:= \lambda \setminus \{c_2, c_3\} = (4, 2, 1^{q-1}),
      &|\mu_B| = q+5 \\
   \mu_C &:= \lambda \setminus \{c_1, c_2\} = (3, 2, 1^q),
      &|\mu_C| = q+5
\end{align*}

\paragraph{Outer chain.}
$\mathcal{O}_\lambda := (T_{j_0-1}{+}1)(T_{j_0}{+}1)(T_{j_0+1}{+}1)(T_{n-2}{+}1)
= (T_{q+2}{+}1)(T_{q+3}{+}1)(T_{q+4}{+}1)(T_{q+5}{+}1)$,
four factors with indices $\ell, \ell+1, n-3, n-2$.

\paragraph{Recall: May-13-UB reduction.}
\begin{equation}\label{eq:UB-reduction}
   B^{(\lambda)}_{j_0} V_\lambda
   \;=\;
   \mathcal{O}_\lambda \Big[
      \Phi_A^\lambda\big(B^{(\mu_A)}_{j_0(\mu_A)} V_{\mu_A}\big)
      +
      \Phi_B^\lambda\big(B^{(\mu_B)}_{j_0(\mu_B)} V_{\mu_B}\big)
   \Big].
\end{equation}
The $\mu_C$-piece and same-row $S$-piece both vanish structurally for
all $q \ge 3$ (May-13-UB Lemmas 3.5 and 3.6: $j$-formula leakage
and the same-row-dies framework).

\paragraph{Inner basis inputs.}
\begin{itemize}[leftmargin=2em,topsep=0.2em,itemsep=0.1em]
\item $\dim B^{(\mu_A)}_{j_0(\mu_A)} V_{\mu_A} = 2$. Structural at
$q \ge 5$ (SRD-331 + May-12 evening on the inner inner shape
$\mu_1^A := \mu_A \setminus \{c_1(\mu_A), c_3(\mu_A)\}
= (3, 2, 1^{q-2})$). Computational at $q \in \{3, 4\}$. Basis
$\{u_A^1, u_A^2\}$ explicit (via lifting through the r=1 reduction;
see Section~\ref{sec:muA-basis}).
\item $\dim B^{(\mu_B)}_{j_0(\mu_B)} V_{\mu_B} = 3$. Structural at
$q \ge 5$ (May-12 evening-3 + SRD-421). Computational at
$q \in \{3, 4\}$ inputs. Basis $\{u_B^1, u_B^2, u_B^3\}$ explicit
(May-12-Eve3 \S5--7; see Section~\ref{sec:muB-basis}).
\end{itemize}

\section{Five-vector construction}\label{sec:five-vec}

\begin{definition}\label{def:five-vec}
Define five vectors in $V_\lambda$:
\begin{align*}
   F_A^i \;&:=\; \mathcal{O}_\lambda\, \Phi_A^\lambda(u_A^i),
   \qquad i \in \{1, 2\}, \\
   F_B^j \;&:=\; \mathcal{O}_\lambda\, \Phi_B^\lambda(u_B^j),
   \qquad j \in \{1, 2, 3\}.
\end{align*}
\end{definition}

By~\eqref{eq:UB-reduction} together with the inner-basis spanning
properties, $\{F_A^i\}_i \cup \{F_B^j\}_j$ spans
$B^{(\lambda)}_{j_0} V_\lambda$. The conjecture is that these five
vectors are \emph{linearly independent}, which would give
$\dim = 5$.

\begin{theorem}[A--B decoupling reduces lower bound to two within-class projections]\label{thm:main}
Suppose
\begin{enumerate}[leftmargin=2em,topsep=0.2em,itemsep=0.1em,label=(\alph*)]
\item $\{\pi^{(c_1)}_n F_A^1,\ \pi^{(c_1)}_n F_A^2\}$ are
linearly independent in $\pi^{(c_1)}_n V_\lambda$.
\item $\{\pi^{(c_2)}_n F_B^1,\ \pi^{(c_2)}_n F_B^2,\ \pi^{(c_2)}_n F_B^3\}$
are linearly independent in $\pi^{(c_2)}_n V_\lambda$.
\end{enumerate}
Then the five vectors $\{F_A^1, F_A^2, F_B^1, F_B^2, F_B^3\}$ are
linearly independent in $V_\lambda$. In particular,
$\dim B^{(\lambda)}_{j_0} V_\lambda \ge 5$.
\end{theorem}

The structural content of Theorem~\ref{thm:main} is the reduction:
the full $5$-vector independence follows from \emph{two independent}
within-class projection independences. The A vs B coupling is broken
by the position-of-$n$ projection alone.

\section{A vs B separation via position-of-$n$}\label{sec:AvsB}

\subsection{Position-of-$n$ is a $\mathcal{O}_\lambda$-invariant}

\begin{lemma}[Position-of-$n$ invariant under outer chain]\label{lem:posn}
For every seminormal basis vector $v_T \in V_\lambda$ and every
$1 \le j \le n - 2$, $T_j v_T$ is a $\mathbb{Q}(q)$-linear
combination of seminormal basis vectors $v_{T'}$ with
$\mathrm{pos}_{T'}(n) = \mathrm{pos}_T(n)$.
\end{lemma}

\begin{proof}
$T_j$ acts on the seminormal basis via the standard Hoefsmit formulae
(same row $\Rightarrow q$, same column $\Rightarrow -1$, otherwise a
$2$-block on $\{v_T, v_{s_jT}\}$ with $s_j$ swapping letters $j$
and $j+1$). The position of letter $n$ is unaffected by $s_j$ since
$j + 1 \le n - 1$, so neither swap touches the letter~$n$.\qed
\end{proof}

\begin{corollary}[Seminormal projection commutes with $\mathcal{O}_\lambda$]\label{cor:proj-commute}
For each cell $c$ of $\lambda$, the projection $\pi^{(c)}_n$ onto
seminormal vectors with $\mathrm{pos}(n) = c$ commutes with each
factor $(T_j{+}1)$, $1 \le j \le n - 2$. Since the indices in
$\mathcal{O}_\lambda$ are $\{q+2, q+3, q+4, q+5\} \subset
\{1, \ldots, n-2\}$ (as $n - 2 = q + 5$), $\pi^{(c)}_n$ commutes
with $\mathcal{O}_\lambda$.
\end{corollary}

\subsection{Image of $\Phi_X^\lambda$ has restricted position-of-$n$ support}

\begin{lemma}[Position-of-$n$ support of $\Phi_X^\lambda(V_{\mu_X})$]\label{lem:Phi-pos-supp}
For each corner pair $X \in \{A, B, C\}$ of $\lambda$ and every
$v \in V_{\mu_X}$, every SYT $T$ in the seminormal support of
$\Phi_X^\lambda(v)$ has $\mathrm{pos}_T(n) \in \{c_X^-, c_X^+\}$.
In particular:
\begin{itemize}[leftmargin=2em,topsep=0.2em,itemsep=0.1em]
\item For $X = A$: $\mathrm{pos}_T(n) \in \{c_1, c_3\} = \{(1, 4), (q+2, 1)\}$.
\item For $X = B$: $\mathrm{pos}_T(n) \in \{c_2, c_3\} = \{(2, 3), (q+2, 1)\}$.
\end{itemize}
\end{lemma}

\begin{proof}
By the construction of $\Phi_X^\lambda$ (May-13-UB
Theorem 2.2, item 1): $\Phi_X^\lambda(v_{T'}^{(\mu_X)})
= v_{T_A^X} + \tau_X v_{T_B^X}$ where $T_A^X, T_B^X$ are the two
extensions of $T'$ obtained by placing $\{n-1, n\}$ at the corner
pair $\{c_X^-, c_X^+\}$. Both extensions have
$\mathrm{pos}(n) \in \{c_X^-, c_X^+\}$. Extending by linearity to
$V_{\mu_X}$, the seminormal support of $\Phi_X^\lambda(v)$ consists
of such extensions, hence preserves the position-of-$n$ support
restriction.\qed
\end{proof}

\subsection{Cross-vanishings}

\begin{proposition}[Cross-projection vanishings at $\lambda$]\label{prop:cross-zero}
For all $i \in \{1, 2\}$ and $j \in \{1, 2, 3\}$,
\[
   \pi^{(c_2)}_n F_A^i \;=\; 0,
   \qquad
   \pi^{(c_1)}_n F_B^j \;=\; 0.
\]
\end{proposition}

\begin{proof}
By Corollary~\ref{cor:proj-commute}, $\pi^{(c_2)}_n F_A^i
= \pi^{(c_2)}_n \mathcal{O}_\lambda \Phi_A^\lambda(u_A^i)
= \mathcal{O}_\lambda \pi^{(c_2)}_n \Phi_A^\lambda(u_A^i)$. By
Lemma~\ref{lem:Phi-pos-supp}, every SYT in the seminormal support of
$\Phi_A^\lambda(u_A^i)$ has $\mathrm{pos}(n) \in \{c_1, c_3\}$,
\emph{never} $c_2$. Hence $\pi^{(c_2)}_n \Phi_A^\lambda(u_A^i) = 0$,
so $\pi^{(c_2)}_n F_A^i = 0$. Symmetrically for the
$F_B^j$ case.\qed
\end{proof}

\subsection{Conclusion: $A$ and $B$ are separated}

\begin{proof}[Proof of Theorem~\ref{thm:main}]
Suppose
\begin{equation}\label{eq:vanishing}
   \alpha_1 F_A^1 + \alpha_2 F_A^2 + \beta_1 F_B^1
   + \beta_2 F_B^2 + \beta_3 F_B^3 \;=\; 0.
\end{equation}
Apply $\pi^{(c_2)}_n$ to both sides. By
Proposition~\ref{prop:cross-zero}, $\pi^{(c_2)}_n F_A^i = 0$ for both
$i$. Hence
\[
   \beta_1 \pi^{(c_2)}_n F_B^1
   + \beta_2 \pi^{(c_2)}_n F_B^2
   + \beta_3 \pi^{(c_2)}_n F_B^3 \;=\; 0.
\]
By hypothesis (b) of Theorem~\ref{thm:main}, the vectors
$\{\pi^{(c_2)}_n F_B^j\}_{j=1}^3$ are linearly independent in
$\pi^{(c_2)}_n V_\lambda$. Therefore $\beta_1 = \beta_2 = \beta_3 = 0$.

Substituting back into~\eqref{eq:vanishing}:
$\alpha_1 F_A^1 + \alpha_2 F_A^2 = 0$. Apply $\pi^{(c_1)}_n$ to this
equation: $\alpha_1 \pi^{(c_1)}_n F_A^1 + \alpha_2 \pi^{(c_1)}_n F_A^2 = 0$.
By hypothesis (a), $\{\pi^{(c_1)}_n F_A^i\}_{i=1}^2$ are linearly
independent. So $\alpha_1 = \alpha_2 = 0$.

Hence all five coefficients vanish, proving linear independence.
Since these five vectors lie in $B^{(\lambda)}_{j_0} V_\lambda$
(by~\eqref{eq:UB-reduction} and Definition~\ref{def:five-vec}), we
have $\dim B^{(\lambda)}_{j_0} V_\lambda \ge 5$.\qed
\end{proof}

\begin{remark}[No $c_3$-analysis needed]\label{rem:no-c3}
The proof does \emph{not} require any analysis of the
$\pi^{(c_3)}_n$-components: the position-of-$n$ projection alone
suffices to decouple $A$ from $B$, provided the within-class
projection independences (a), (b) are established. The
$c_3$-direction, where both $\Phi_A^\lambda$- and
$\Phi_B^\lambda$-supports overlap, plays no role.
\end{remark}

\begin{proposition}[Discarded sub-claim: NOT TRUE in general]\label{prop:discard}
The naive sub-statement ``\,$\alpha_1 F_A^1 + \alpha_2 F_A^2 = 0$ and
$\beta_1 F_B^1 + \beta_2 F_B^2 + \beta_3 F_B^3 = 0$ separately,
from~\eqref{eq:vanishing} alone'' is \emph{stronger than}
Theorem~\ref{thm:main} and is not directly provable from the
position-of-$n$ argument alone (both sides could lie non-trivially
in the shared $\pi^{(c_3)}_n V_\lambda$). Theorem~\ref{thm:main}'s
statement (via the hypothesis (a), (b)) is the correct, sufficient
formulation.
\end{proposition}


\section{Inner-basis structure}\label{sec:inner-basis}

\subsection{The $\mu_A$ basis via the r=1 reduction}\label{sec:muA-basis}

$\mu_A = (3, 3, 1^{q-1})$ has $\ell(\mu_A) = q+1$, $j_0(\mu_A) = q+2$,
$|\mu_A| = q+5$. It has a unique length-preserving corner
$c_1(\mu_A) = (2, 3)$, column-$1$ bottom $c_*(\mu_A) = (q+1, 1)$, and
a same-row pair at row $2$.

The r=1 reduction (May-13 r1 + SRD-331) reads
\begin{equation}\label{eq:muA-reduction}
   B^{(\mu_A)}_{j_0(\mu_A)} V_{\mu_A}
   \;=\;
   \mathcal{O}_{\mu_A}\, \Phi_*^{\mu_A}\big(B^{(\mu_1^A)}_{j_0(\mu_1^A)} V_{\mu_1^A}\big),
\end{equation}
where $\mu_1^A := \mu_A \setminus \{c_1(\mu_A), c_*(\mu_A)\}
= (3, 2, 1^{q-2})$ on $q+3$ letters, and
$\mathcal{O}_{\mu_A} := (T_{q+1}{+}1)(T_{q+2}{+}1)(T_{q+3}{+}1)$
(three factors).

$\dim B^{(\mu_1^A)}_{j_0(\mu_1^A)} V_{\mu_1^A} = 2$ by the
May-12 \texttt{2026-05-12-dim-2-321.tex} paper (structural at
$|\mu_1^A| \ge 8$, i.e., $q \ge 5$; computational at $q \in \{3, 4\}$).
Explicit basis
\[
   \widetilde F_A^{\mu_1^A} := \mathcal{O}_{\mu_1^A}\,
        \Phi_A^{\mu_1^A}\big(u^{(\nu_A^{\mu_1^A})}\big),
   \qquad
   \widetilde F_B^{\mu_1^A} := \mathcal{O}_{\mu_1^A}\,
        \Phi_B^{\mu_1^A}\big(u^{(\nu_B^{\mu_1^A})}\big),
\]
where $\nu_A^{\mu_1^A} = (2, 2, 1^{q-3})$ and $\nu_B^{\mu_1^A} = (3, 1^{q-2})$
on $q+1$ letters, $\mathcal{O}_{\mu_1^A} = (T_q{+}1)(T_{q+1}{+}1)$
(two factors), and $u^{(\nu)}$ denotes the dim-$1$ generator of the
$j$-formula image on $\nu$.

\begin{definition}[$\mu_A$-basis]\label{def:muA-basis}
\[
   u_A^1 := \mathcal{O}_{\mu_A}\, \Phi_*^{\mu_A}\big(\widetilde F_A^{\mu_1^A}\big),
   \qquad
   u_A^2 := \mathcal{O}_{\mu_A}\, \Phi_*^{\mu_A}\big(\widetilde F_B^{\mu_1^A}\big).
\]
\end{definition}

By~\eqref{eq:muA-reduction}, $\{u_A^1, u_A^2\}$ spans $B^{(\mu_A)}$.
The lower bound $\dim \ge 2$ requires showing
$\mathcal{O}_{\mu_A} \Phi_*^{\mu_A}$ is injective on the
$2$-dimensional $B^{(\mu_1^A)}$; this is an open structural problem
(verified computationally at $q \in \{3, 4, 5\}$, the SRD-331 paper's
range).

\subsection{The $\mu_B$ basis from May-12-Eve3}\label{sec:muB-basis}

$\mu_B = (4, 2, 1^{q-1})$ is the May-12 evening-3 family at
$|\mu_B| = q+5$ (with the paper's $n = q+5$). The three-vector basis
\[
   u_B^1 := F_{A,A}^{(\mu_B)}, \quad
   u_B^2 := F_{A,B}^{(\mu_B)}, \quad
   u_B^3 := F_B^{(\mu_B)}
\]
is May-12-Eve3 Definition~5.1, with linear independence proved there
(May-12-Eve3 \S5.3, \S6, \S7; \texttt{2026-05-12-evening-4-2-1n-dim3.tex}).

\section{Why naive position-of-letter separators fail}\label{sec:obstruction}

\subsection{The natural candidate: position-of-$(n-2)$ within $\pi^{(c_1)}_n$}

For the $(4, 3, 1^q)$ case, the within-A separation problem is:
distinguish $\pi^{(c_1)}_n F_A^1$ from $\pi^{(c_1)}_n F_A^2$.

Both are seminormal vectors in $\pi^{(c_1)}_n V_\lambda$ (the subspace
of $V_\lambda$ supported on SYTs with letter $n = q+6$ at
$c_1 = (1, 4)$). Within this subspace, $n - 1 = q+5$ has position
$c_3 = (q+2, 1)$ for all vectors in the support of
$\pi^{(c_1)}_n \Phi_A^\lambda(\cdot)$ (the ``B-direction'' of
$\Phi_A^\lambda$).

The natural separator candidate: position-of-$(n-2) = q+4$.

\begin{lemma}[Position-of-$(n-2)$ is moved by both $(T_{n-3}+1)$ and $(T_{n-2}+1)$]\label{lem:posn-2-move}
$T_j$ preserves position-of-$(n-2)$ iff $j \notin \{n-3, n-2\}$.
\end{lemma}

\begin{proof}
$T_j$ acts on the pair $\{j, j+1\}$. Position-of-$k$ is invariant
under $T_j$ iff $\{j, j+1\} \cap \{k\} = \emptyset$, i.e.,
$j \ne k$ and $j \ne k-1$.\qed
\end{proof}

Inside $\mathcal{O}_\lambda = (T_{q+2}+1)(T_{q+3}+1)(T_{q+4}+1)(T_{q+5}+1)$,
the indices are $\{q+2, q+3, q+4, q+5\}$. Position-of-$(n-2) = q+4$ is
moved by $T_{q+3}$ (which acts on $\{q+3, q+4\}$) and by $T_{q+4}$
(which acts on $\{q+4, q+5\}$). The factors $(T_{q+2}+1)$ and
$(T_{q+5}+1)$ preserve it (the former touches $\{q+2, q+3\}$ only,
the latter $\{q+5, q+6\}$).

\subsection{Two-step contamination}

The contamination pattern is structural and worth stating in general
terms. Consider the recursive lift through three nested partitions:
\[
   \mu_1^A = (3, 2, 1^{q-2}) \;\subset\; \mu_A = (3, 3, 1^{q-1}) \;\subset\; \lambda = (4, 3, 1^q),
\]
with intermediate sizes $q+3, q+5, q+7$ respectively.

At each layer, the basis vectors $\widetilde F^{(\mu^{\text{inner}})}_A$
and $\widetilde F^{(\mu^{\text{inner}})}_B$ at the inner level are
distinguished by the position of the inner top letter (call it
$k^{\text{inner}}$) at some pair of exclusive cells.

\begin{lemma}[Position-of-letter moves under outer chain]\label{lem:pos-move}
For a chain
$\mathcal{O} = (T_{k_1}+1)(T_{k_2}+1)\cdots(T_{k_r}+1)$, the position
of letter $\ell$ is preserved by $T_{k_i}$ iff
$\ell \notin \{k_i, k_i + 1\}$. Hence the cells reachable by the
position-of-$\ell$ orbit after $\mathcal{O}$ form an ``orbit'' under
the chain of moves
$\{T_{k_i} : k_i = \ell - 1 \text{ or } \ell\}$.
\end{lemma}

\begin{proposition}[Contamination at the $\lambda$ level]\label{prop:contamination}
For the natural inner separator at every layer, the outer chain
$\mathcal{O}_\lambda$ at the $\lambda$ level enlarges the
position-of-letter orbit to overlap with the orbit of the other inner
basis vector. Consequently, no single position-of-letter projection
$\pi^{(c)}_k$ at the $\lambda$ level is exclusive between
$F_A^1$ and $F_A^2$ (or between any $F_X^i$ pair sharing the same
$\Phi_X^\lambda$ class).
\end{proposition}

\begin{proof}[Proof sketch]
The position of any specific letter $k$ at the $\lambda$ level is
determined by:
\begin{enumerate}[leftmargin=2em,topsep=0.2em,itemsep=0.1em]
\item Its position in the original $\mu_1^A$- (or $\mu_B$-)level basis
vector.
\item Possible moves through $\mathcal{O}_{\mu_1^A}$ at the innermost
layer (two $(T_j+1)$ factors).
\item Lift through $\Phi_*^{\mu_A}$ (or analog).
\item Moves through $\mathcal{O}_{\mu_A}$ at the middle layer (three
factors).
\item Lift through $\Phi_X^\lambda$.
\item Moves through $\mathcal{O}_\lambda$ at the outer layer (four
factors).
\end{enumerate}
Total: $2 + 3 + 4 = 9$ $(T_j+1)$ factors in the cumulative chain.
Each $T_j$-move can replace position-of-letter-$k$ with the current
position of letter $k-1$ or $k+1$ (depending on the index match
with the move).

Starting from any exclusive cell at the innermost layer, the
$9$-step orbit reaches a set of cells that, for typical Hoefsmit
non-degenerate configurations, includes the originally-exclusive
cell of the OTHER inner basis vector. Concretely: the cells
$(1, 3)$ and $(2, 2)$ are at axial distance $c(1, 3) - c(2, 2) = 2$,
which yields a non-degenerate $T_j$-$2$-block for appropriate $j$, and
this move \emph{is} reachable somewhere in the $9$-step chain. Hence
the contamination occurs generically.

A fully rigorous proof would enumerate the $9$-step orbit
explicitly (a finite case analysis, parallel to May-12-Eve3 \S5.4 but
expanded by an order of magnitude due to the additional recursion
depth). We omit the enumeration; the structural conclusion suffices
for the discussion.
\qed
\end{proof}

\begin{remark}[Cost of recursion depth]\label{rem:depth-cost}
The recursion depth of the contamination is structurally tied to the
shape:
\begin{itemize}[leftmargin=2em,topsep=0.2em,itemsep=0.1em]
\item Hooks (e.g.\ $(k, 1^*)$): one layer of recursion; clean
position-of-$(n-1)$ separators (May-12 hook paper).
\item $r = 2$ shapes without long inner non-hook chains (e.g.\
$(4, 2, 1^*)$): two layers; position-of-$(n-1)$ + position-of-$(n-2)$
suffice (May-12-Eve3).
\item $r = 2$ same-row shapes like $(4, 3, 1^*)$: three layers (the
$\mu_A$-side has a same-row pair triggering an extra r=1 lift). The
$9$-step chain contaminates the natural separator.
\end{itemize}
This suggests a coarse rule: \emph{position-of-letter techniques
exhaust at depth $\le 2$.} Deeper recursions require finer tools.
\end{remark}

\section{Path forward: tracer-pair technique}\label{sec:tracer}

The memory note \texttt{tracer-pair-technique} (this session) records
a successful method from the hook program:

\begin{quote}\emph{Tracer-pair:} when $\dim \le d$ is known, track $d$
matrix coefficients between explicit SYTs; fixed entries force unique
contributors.\end{quote}

\subsection{Applied to $(4, 3, 1^q)$ lower bound}

Choose explicit $\lambda$-SYTs $T^*_1, \ldots, T^*_5$ and compute
the $5 \times 5$ ``tracer matrix''
\[
   M_{ik} := \langle v_{T^*_i}, F_X^{j_k}\rangle
\]
where $(X, j_k)$ runs through the $5$ basis indices
$\{(A, 1), (A, 2), (B, 1), (B, 2), (B, 3)\}$. If $\det M \in
\mathbb{Q}(q)^\times$, then $F$'s are linearly independent.

\subsection{Reduction via the A--B decoupling (Theorem~\ref{thm:main})}

By Theorem~\ref{thm:main} (the A vs B separation), the
$5 \times 5$ tracer matrix has block-diagonal structure: choose two
$T^*_i$'s with $\mathrm{pos}(n) = c_1$ (capturing only $F_A$'s) and
three with $\mathrm{pos}(n) = c_2$ (capturing only $F_B$'s). Then $M$
splits into a $2 \times 2$ block (for $F_A$'s) and a $3 \times 3$
block (for $F_B$'s).

\subsection{The two reduced problems}

\paragraph{$2 \times 2$ tracer for $F_A^1, F_A^2$.}
Pick two $\lambda$-SYTs $T^*_1, T^*_2$ with $\mathrm{pos}(n) = c_1$.
Distinguish them by some interior letter pattern that detects the
inner $\widetilde F_A^{\mu_1^A}$ vs $\widetilde F_B^{\mu_1^A}$
distinction. Concretely: $T^*_1$ has interior column-$2$ entry
$\gamma \in \{q, q+1, q+2\}$ (from the $u^{\nu_A}$ support, before
all the chains); $T^*_2$ has interior row-$1$ entry $b = q+2$
(distinguishing the $\widetilde F_B^{\mu_1^A}$-derived inner basis).

\paragraph{$3 \times 3$ tracer for $F_B^1, F_B^2, F_B^3$.}
Pick three $\lambda$-SYTs $T^*_3, T^*_4, T^*_5$ with
$\mathrm{pos}(n) = c_2$. The May-12-Eve3 separation (column-$2$ entries
of the inner $\mu_A^{\mu_B} = (3, 2, 1^{q-2})$ shape vs row-$1$ entries
of the inner $\mu_B^{\mu_B} = (4, 1^{q-1})$ shape) provides three
clean SYT supports.

\subsection{Status}

The tracer-pair technique reduces the structural lower bound to
\emph{five explicit Hoefsmit coefficient computations} (the
$5 \times 5 = 4 + 9$ entries of the $2 \times 2$ and $3 \times 3$
blocks). Each entry is a product of:
\begin{enumerate}[leftmargin=2em,topsep=0.2em,itemsep=0.1em]
\item One ``inner-inner'' Hoefsmit coefficient (from
$\mathcal{O}_{\mu_1^A}$ acting on $u^\nu$).
\item One ``inner'' Hoefsmit coefficient (from $\mathcal{O}_{\mu_A}$
or $\mathcal{O}_{\mu_B}$).
\item One ``outer'' Hoefsmit coefficient (from $\mathcal{O}_\lambda$).
\item Several $\tau, \tau_X$ factors from the $\Phi$-maps.
\end{enumerate}
The non-vanishing of the $2 \times 2$ and $3 \times 3$ determinants
reduces to algebraic non-vanishing of explicit polynomial
expressions in $q$. This is a finite computation, verifiable at any
specific $q$ and (with care) provable structurally via the $b' = 1$
propagation technique (memory note
\texttt{b-prime-propagation}).

\section{Computational verification}\label{sec:verify}

From the May-13-UB paper Section~5 (verbatim):
\begin{center}
\begin{tabular}{lcc}
\toprule
$\lambda$ & $\dim V_\lambda$ & $\dim B^{(\lambda)}_{j_0} V_\lambda$ \\
\midrule
$(4, 3, 1, 1, 1)$ \quad ($q = 3$, $n = 10$) & $525$ & $5$ \\
$(4, 3, 1, 1, 1, 1)$ \quad ($q = 4$, $n = 11$) & $1100$ & $5$ \\
\bottomrule
\end{tabular}
\end{center}
At both verified cases, $\dim = 5$, consistent with the conjectural
lower bound from the five-vector construction.

A separate, more refined computational check (not yet performed in
this session): compute the explicit five vectors $F_A^1, F_A^2, F_B^1,
F_B^2, F_B^3$ via the chain of operations defined in
Section~\ref{sec:five-vec} and verify that the resulting $5$-vector
matrix has rank $5$. This would also exhibit the supports and
enable a precise tracer-pair selection.

\section{Honest status}\label{sec:honest}

\paragraph{What we proved structurally.}
\begin{enumerate}[leftmargin=2em,topsep=0.2em,itemsep=0.1em]
\item The explicit five-vector construction
(Definition~\ref{def:five-vec}).
\item Position-of-$n$ separation between A-pieces and B-pieces
(Proposition~\ref{prop:cross-zero}).
\item The full decoupling statement (Theorem~\ref{thm:main} +
proof in \S\ref{sec:AvsB}): the $5$-vector independence reduces
to two within-class projection independences, without any analysis
of the $\pi^{(c_3)}_n$-component (Remark~\ref{rem:no-c3}).
\end{enumerate}

\paragraph{What we identified.}
\begin{enumerate}[leftmargin=2em,topsep=0.2em,itemsep=0.1em]
\item The naive position-of-letter separator (analogous to
May-12-Eve3 Lemma~5.4) does \emph{not} propagate cleanly through
multi-layer recursions (Proposition~\ref{prop:contamination}). The
$(4, 3, 1^q)$ shape requires three layers of recursion, and the
chain of $T_j$ moves contaminates the natural separator at every
layer.
\item The correct technique is \emph{tracer-pair} via explicit
Hoefsmit coefficient computation (Section~\ref{sec:tracer}). The
reduction splits the $5$-vector problem into a $2 \times 2$ and a
$3 \times 3$ tracer computation.
\end{enumerate}

\paragraph{What remains.}
\begin{enumerate}[leftmargin=2em,topsep=0.2em,itemsep=0.1em]
\item Execute the $2 \times 2$ tracer for $\{F_A^1, F_A^2\}$ — five
nested Hoefsmit factors, structural at $q \ge q_*$, computational at
small $q$. Reduces to the structural lower bound for
$\mu_A = (3, 3, 1^{q-1})$, currently open
(SRD-331 Remark~3.4).
\item Execute the $3 \times 3$ tracer for
$\{F_B^1, F_B^2, F_B^3\}$ — directly lifts the May-12-Eve3
position-of-$(n-1)$ analysis, plus the outer factor $(T_{n-5}+1)$
specific to $(4, 3, 1^q)$.
\end{enumerate}
(The $c_3$-direction analysis is \emph{not} needed —
Remark~\ref{rem:no-c3} — only the two within-class projection
independences above.)

\paragraph{Combined with May-13-UB.}
If the three remaining items are closed, we obtain
$\dim B^{(\lambda)}_{j_0} V_\lambda = 5$ structurally for all
$q \ge q_*$ (with $q_*$ to be determined by the chain of structural
ranges).

\section{Comparison with May-12-Eve3}\label{sec:comparison}

The May-12-Eve3 paper proved $\dim B^{((4, 2, 1^{n-6}))} = 3$ by
the analogous strategy:
\begin{itemize}[leftmargin=2em,topsep=0.2em,itemsep=0.1em]
\item Two-vector A-piece $F_{A,A}, F_{A,B}$ via inner basis from
May-12-today on $\mu_A = (3, 2, 1^{m-5})$.
\item One-vector B-piece $F_B$ via inner basis from May-12 hook on
$\mu_B = (4, 1^{m-4})$.
\item A vs B separation via position-of-$n$.
\item Within-A separation via position-of-$(n-1)$, with a single
$T_{n-2}$ outer-factor move analysed via Lemma~posnm1.
\end{itemize}

The key technical step that worked there: the outer chain had
\emph{three} factors $(T_{n-4}+1)(T_{n-3}+1)(T_{n-2}+1)$, and the
position-of-$(n-1)$ was preserved by the first two factors. The
single $T_{n-2}$-move was sufficient to track and the reachable
cells were a $2$-element set with clean exclusive elements.

For $(4, 3, 1^q)$, the outer chain has \emph{four} factors. The
analogous separator position-of-$(n-2)$ would be preserved by only
two of the four factors, and the two outer-factor moves enlarge the
reachable set beyond the $2$-element ``clean'' structure.

\emph{This is a structural difficulty inherent to the shape, not a
defect of the proof method.} It explains why the May-12-Eve3
technique does not directly translate, and motivates the
tracer-pair refinement.

\end{document}
